\documentclass[12pt]{article}
\usepackage[utf8]{inputenc}
\usepackage[T1]{fontenc}
\usepackage{amsmath,amssymb,amsthm}
\usepackage{booktabs}
\usepackage{float}
\usepackage[margin=1in]{geometry}
\usepackage{setspace}

\usepackage[hidelinks]{hyperref}

\onehalfspacing
\setlength{\emergencystretch}{2em}
\makeatletter\def\verbatim@font{\normalfont\ttfamily\small}\makeatother
\newtheorem{proposition}{Proposition}[section]
\theoremstyle{definition}
\newtheorem{remark}[proposition]{Remark}
\renewcommand{\thesection}{S\arabic{section}}
\renewcommand{\theequation}{S\arabic{equation}}
\renewcommand{\thetable}{S\arabic{table}}
\renewcommand{\theproposition}{S\arabic{section}.\arabic{proposition}}
\newcommand{\ki}{k_i}
\newcommand{\kj}{k_j}
\newcommand{\gt}{\tilde{g}}
\title{Supplementary Material\\
\large Dividing by a difference of noisy baselines:\\ why per-dyad estimates of a coupling parameter fail, and how to size a pooled alternative}
\author{Mark D. Izgarshev and Sophie Noemi Kieckens}
\date{}
\begin{document}
\maketitle

\noindent Equations, sections and propositions prefixed S refer to this document;
unprefixed equation numbers refer to the main text.

% =====================================================================
\section{The coupling models}

\subsection{A shared multiplicative gain cannot produce convergence}

A natural first formulation treats company as raising both diners' engagement
through a factor common to the dyad:
\begin{equation}\label{eq:shared}
\ki'(t)=\beta_ie^{-l_it}\,w(t),
\qquad
\kj'(t)=\beta_je^{-l_jt}\,w(t),
\qquad w>0.
\end{equation}
Then $D'=\kj'-\ki'=(\beta_je^{-l_jt}-\beta_ie^{-l_it})\,w$, so
\[
\operatorname{sgn}D'(t)=\operatorname{sgn}\bigl(\beta_je^{-l_jt}-\beta_ie^{-l_it}\bigr)
\]
for any positive $w$. Nothing in the argument uses more than positivity, so $w$ may
depend on time, on the state, or on both. For $n\ge2$ agents with
$x_a'=f_a(t)\,w(t,x_1,\ldots,x_n)$ and $w>0$,
$\operatorname{sgn}(x_a-x_b)'=\operatorname{sgn}(f_a-f_b)$ for every pair.

This fixes the \emph{direction} in which each gap moves, not whether it shrinks.
If $\int^\infty w<\infty$ the differences freeze at nonzero values. When a gap does
contract, the contraction comes from the individual profiles, since $|x_a-x_b|$
decreases exactly when $x_a-x_b$ and $f_a-f_b$ have opposite signs. The form
contains no restoring term whose sign responds to the current gap.

\paragraph{An exactly solvable case.} Take $w=a(1+|\ki-\kj|)$, with $a>0$ a
closeness score. Then $D'=g(1+|D|)$ with
$g=a(\beta_je^{-l_jt}-\beta_ie^{-l_it})$. Let
$\psi(x)=\operatorname{sgn}(x)\log(1+|x|)$, so $\psi'(x)=1/(1+|x|)$ and
$\psi(D(t))=\psi(D(0))+\Phi(t)$ with $\Phi(t)=\int_0^tg$. For a shared meal,
$D(0)=0$ and $D(t)=\operatorname{sgn}\Phi(t)\,(e^{|\Phi(t)|}-1)$. As $t\to\infty$,
$\Phi\to\Phi_\infty=a(\beta_j/l_j-\beta_i/l_i)$ and
\[
|D_\infty|=e^{|\Phi_\infty|}-1,
\]
which grows without bound in closeness: the closer two diners feel, the further
apart they finish. This is the formulation of Kieckens and Izgarshev (2026), who
coupled the diners through a shared affiliation factor
$a_{ij}(|\ki-\kj|+1)$, with $a_{ij}$ a closeness score, and who note that the factor
cannot steer the gap. The argument above is more general, since it holds for any
positive shared factor, and it motivates the antisymmetric couplings of the main
text.

\subsection{Intake matching}

For main-text Eq.~(1) with $D(0)=D_0$, the gap obeys $D'=\gt-2\kappa D$,
which is linear. Multiplying by $e^{2\kappa t}$ gives
$(De^{2\kappa t})'=\gt e^{2\kappa t}$, so
\begin{equation}\label{eq:voc}
D(t)=D_0e^{-2\kappa t}+\int_0^t e^{-2\kappa(t-s)}\gt(s)\,ds .
\end{equation}
Using $\int_0^t\beta e^{-ls}e^{2\kappa s}ds=\beta(e^{(2\kappa-l)t}-1)/(2\kappa-l)$,
for $2\kappa\neq l_i,l_j$,
\begin{equation}\label{eq:solution}
D(t)=D_0 e^{-2\kappa t}
+\frac{A\beta_j(e^{-l_j t}-e^{-2\kappa t})}{2\kappa-l_j}
-\frac{A\beta_i(e^{-l_i t}-e^{-2\kappa t})}{2\kappa-l_i},
\end{equation}
with the removable limit $A\beta\,te^{-lt}$ in the excluded cases. For $\kappa>0$,
$D(t)\to0$ at the exponential rate $\min(2\kappa,l_i,l_j)$, since every term is a
multiple of $e^{-l_it}$, $e^{-l_jt}$ or $e^{-2\kappa t}$. That rate is informative
about the coupling only when $2\kappa<\min(l_i,l_j)$; otherwise the approach to zero
is set by the decaying drives and would occur with $\kappa=0$. The convergence is
also asymptotic in a system whose forcing decays over hours, so it is not a
prediction about meals. The main text therefore works on $[0,T]$.

Writing $\Sigma(T)=\ki(T)+\kj(T)$ for the conserved total of Eq.~(2), integration
gives $\ki(T)=\tfrac12(\Sigma(T)-D(T))$ and $\kj(T)=\tfrac12(\Sigma(T)+D(T))$. Each
diner is at the mean of the two solo intakes, offset by half the residual gap; the
offset vanishes only as $\kappa\to\infty$ or $T\to\infty$.

\paragraph{Positivity.} Because the coupling subtracts from the leading diner's
rate, a bite frequency can in principle become negative. Suppose $\ki>\kj$, so
$D<0$. From \eqref{eq:voc} with $D_0=0$ and
$\gt(s)\geq-A\beta_ie^{-l_is}$,
\[
\kappa|D(t)|\le \kappa A\beta_i\frac{e^{-l_it}-e^{-2\kappa t}}{2\kappa-l_i}
\le A\beta_ie^{-l_it}\,\frac{\kappa}{2\kappa-l_i}\qquad(2\kappa>l_i),
\]
which is at most $A\beta_ie^{-l_it}$ when $\kappa\ge l_i$. Hence $\ki'\ge0$, and by
symmetry $\kj'\ge0$, whenever $\kappa\ge\max(l_i,l_j)$. This is a sufficient
condition. It holds automatically in the slow-drive regime, but it excludes the
regime $2\kappa<\min(l_i,l_j)$ in which the asymptotic convergence rate would be
informative about $\kappa$.

\subsection{Rate matching}

Coupling the rates through their instantaneous difference with a dimensionless
constant $c$, $\ki'=A\beta_ie^{-l_it}+c(\kj'-\ki')$ and symmetrically, is not a
dynamical model. Writing $R=\kj'-\ki'$ gives $R=\gt-2cR$, an algebraic relation
$R=\gt/(1+2c)$ with no memory. The dynamical formulation is Eq.~(3), which makes
the rates state variables.

The matrix of Eq.~(3) has trace $-(l_i+l_j+2\kappa)<0$, determinant
$\Delta:=l_il_j+\kappa(l_i+l_j)>0$ and discriminant $(l_i-l_j)^2+4\kappa^2\ge0$, so
its eigenvalues
\[
\lambda_\pm=-\tfrac12\Bigl[(l_i+l_j+2\kappa)\mp\sqrt{(l_i-l_j)^2+4\kappa^2}\Bigr]
\]
are real and strictly negative for positive $l_i,l_j$, and $r_i,r_j\to0$ for every
$\kappa\ge0$. The closed positive quadrant is forward invariant: if $r_i=0$ and
$r_j\ge0$ then $r_i'=\kappa r_j\ge0$, so the vector field does not point out of the
quadrant on its boundary.

Let $I_a=\int_0^\infty r_a$, finite because the eigenvalues are negative.
Integrating Eq.~(3) over $[0,\infty)$ gives
$(l_i+\kappa)I_i-\kappa I_j=A\beta_i$ and $-\kappa I_i+(l_j+\kappa)I_j=A\beta_j$,
with determinant $\Delta$. Solving, with $D(0)=0$,
$D_\infty(\kappa)=A(\beta_jl_i-\beta_il_j)/\Delta$ and
\begin{equation}\label{eq:ratecontraction}
\frac{D_\infty(\kappa)}{D_\infty(0)}=\frac{1}{1+\kappa(l_i^{-1}+l_j^{-1})}.
\end{equation}
This depends on $l$ through $l^{-1}$ and inherits the non-identifiability of $l$
at meal length; the estimators use the finite-horizon limit below, not this
expression. Unlike intake matching, rate matching does not conserve total intake:
$I_i+I_j=A[\beta_il_j+\beta_jl_i+2\kappa(\beta_i+\beta_j)]/\Delta$, which equals
$A(\beta_i/l_i+\beta_j/l_j)$ at $\kappa=0$ and tends to
$2A(\beta_i+\beta_j)/(l_i+l_j)$ as $\kappa\to\infty$. It is independent of $\kappa$
exactly when $l_i=l_j$. In that case
$r_j-r_i=A(\beta_j-\beta_i)e^{-(l+2\kappa)t}$ and
$D(t)=\frac{A(\beta_j-\beta_i)}{l+2\kappa}\bigl(1-e^{-(l+2\kappa)t}\bigr)$.

\subsection{The slow-drive limit}

\begin{proposition}\label{prop:flat}
Fix $T>0$, $\kappa\neq0$ and $D(0)=0$, and let $l_i,l_j\to0$ with
$\beta_i,\beta_j$ fixed. For both couplings,
\[
\rho(T;\kappa)\;\longrightarrow\;h(\theta)=\frac{1-e^{-\theta}}{\theta},
\qquad\theta=2\kappa T,
\]
independently of the rate at which $l_i,l_j$ vanish and of $\beta_i,\beta_j$.
\end{proposition}
\begin{proof}
For intake matching, $\gt(t)\to A(\beta_j-\beta_i)=:\gamma$ uniformly on $[0,T]$,
so by \eqref{eq:voc}
$D(T)\to\gamma\int_0^Te^{-2\kappa(T-s)}\,ds=\gamma(1-e^{-2\kappa T})/2\kappa$, while
the uncoupled gap tends to $\gamma T$. For rate matching with $l_i=l_j=l$, the
exact ratio is
$\rho=\frac{l}{l+2\kappa}\cdot\frac{1-e^{-(l+2\kappa)T}}{1-e^{-lT}}$, and
$1-e^{-lT}=lT+O((lT)^2)$ gives the same limit; unequal $l$ adds terms of order
$(l_i-l_j)T$.
\end{proof}

$h$ is smooth on $\mathbb{R}$ with $h(0)=1$ (by continuity), and
$h'(\theta)=\theta^{-2}[(1+\theta)e^{-\theta}-1]<0$ for $\theta\ne0$, since
$(1+\theta)e^{-\theta}<1$ for all $\theta\neq0$; $h'(0)=-\tfrac12$. Hence $h$ is
strictly decreasing on the whole real line, with $h>1$ for $\theta<0$ (divergence),
$h\to\infty$ as $\theta\to-\infty$ and $h(\theta)\sim\theta^{-1}$ as
$\theta\to\infty$. Every positive value of $h$ is attained exactly once.

\paragraph{The two couplings are not separable in the mean at meal scale.} They
share the slow-drive contraction factor. When $l_i=l_j$ they also share the total,
since $(r_i+r_j)'=-l(r_i+r_j)$ under Eq.~(3) and the coupling cancels in Eq.~(1).
The leading difference in the mean is $O((l_i-l_j)T)$. They may differ in the
second moment of the gap (Section~\ref{sec:gap}).

\subsection{The detectability ceiling}

Let $\kappa=c\bar s$ with $\bar s\in[s_{\min},s_{\max}]$ and $q=s_{\max}/s_{\min}$,
and let $\Sigma$ be the spread in the contraction factor across the scale. For the
asymptotic factor $\rho=(1+\mu)^{-1}$, with $\mu=\kappa/l_h$ and
$l_h^{-1}=l_i^{-1}+l_j^{-1}$,
\[
\Sigma(\mu)=\frac{1}{1+\mu}-\frac{1}{1+q\mu}=\frac{(q-1)\mu}{(1+\mu)(1+q\mu)} .
\]
Setting $\Sigma'=0$ gives $q(1+\mu)^2=(1+q\mu)^2$, hence $\sqrt q(1+\mu)=1+q\mu$ and
$\mu_{\min}=q^{-1/2}$, at which $\Sigma^\star=(\sqrt q-1)/(\sqrt q+1)$. $\Sigma$ is
unimodal, since $\Sigma(0)=\Sigma(\infty)=0$ and the stationary point is unique.
For $h$ the optimum is at the unique positive root of $h'(\theta)=q\,h'(q\theta)$,
found numerically.

For $q=7$, $\Sigma^\star=0.451$ at $\mu_{\min}=0.378$, with $\rho$ running from
0.726 to 0.274; for $h$, $\theta^\star=0.636$ and $\Sigma^\star=0.518$. Requiring at
least half the maximal spread bounds the position of the bottom of the scale at
$\theta_{\min}\in[0.116,3.14]$, a factor of 27. Main-text Table~1 gives the
degradation under a baseline coupling.

% =====================================================================
\section{The per-dyad estimator}\label{sec:perdyad}

\begin{proposition}\label{prop:estimator}
In the slow-drive limit with $D(0)=0$, $\Gamma=\beta_j-\beta_i$ and
$S=\beta_i+\beta_j$,
\begin{equation}\label{eq:estimator}
h(\theta)=\frac{\kj(T)-\ki(T)}{\ki(T)+\kj(T)}\cdot\frac{S}{\Gamma}.
\end{equation}
\end{proposition}
\begin{proof}
By Proposition~\ref{prop:flat}, $\kj(T)-\ki(T)=A\Gamma T\,h(\theta)$; by Eq.~(2)
in the same limit, $\ki(T)+\kj(T)=AST$. $A$ and $T$ cancel in the ratio.
\end{proof}

\subsection{Precision}

Write $\hat h=a/b$ with $a=\hat S\,(\kj-\ki)/(\kj+\ki)$ and $b=\hat\Gamma$. The
joint session and the solo sessions are independent, so $a$ and $b$ are
independent apart from the weak dependence through $\hat S$. To first order,
$a\approx h\Gamma+\delta_D/(AT)$ with $\mathrm{Var}(\delta_D)=\sigma_D^2$ from main-text
Eq.~(7), plus terms of relative order $(\Gamma/S)^2$ from $\hat S$ and the total,
which are negligible ($(\Gamma/S)^2\approx0.005$ at the FIC median contrast).
Therefore
\begin{equation}\label{eq:var}
\mathrm{Var}(\hat h)\approx\frac{\mathrm{Var}(a)}{\Gamma^2}+\frac{h^2\,\mathrm{Var}(b)}{\Gamma^2}
=\frac{\sigma_D^2}{(A\Gamma T)^2}+\frac{h^2}{\nu^2}=\frac{1}{\zeta^2},
\qquad \nu=\frac{|\Gamma|}{\sigma_{\hat\Gamma}} .
\end{equation}
With $\sigma_o=0$, $\sigma_D^2=\phi AST$ does not depend on $h$, and
$\mathrm{Var}(\hat h)\approx z_0^{-2}+h^2\nu^{-2}$ with $z_0=\Gamma\sqrt{AT/(\phi S)}$.
If the $m$ solo sessions per person have the same length as the joint session and
$A=1$, then $\nu^2=m z_0^2$ and $\mathrm{Var}(\hat h)\approx z_0^{-2}(1+h^2/m)$.

By the delta method, $\mathrm{SD}(\hat\theta)\approx1/(\zeta|h'(\theta)|)$, and for a
normal sampling distribution the IQR is $1.349$ standard deviations. At $\theta=1$,
$|h'(1)|=1-2/e=0.264$, giving $\mathrm{IQR}(\hat\theta)\approx5.11/\zeta$. Simulation
gives $5.23/\zeta$ (range 4.98 to 5.47 across cells with $\zeta>3$); the excess
reflects the curvature of $h^{-1}$. We adopt $\zeta^\star=5.2$, at which the IQR of
$\hat\theta$ equals $\theta=1$.

\paragraph{Required session length.} Requiring $\zeta\ge\zeta^\star$ at
$\theta_{\mathrm{ref}}$, with $h=h(\theta_{\mathrm{ref}})$ and the occasion term of
main-text Eq.~(7),
\begin{equation}\label{eq:Treq}
T_{\mathrm{req}}=\frac{\phi S}{A\Gamma^2}\left[\frac{1}{\zeta^{\star2}}
-\frac{h^2\sigma_{\hat\Gamma}^2}{\Gamma^2}-\frac{2h^2\sigma_o^2}{\Gamma^2}\right]^{-1},
\end{equation}
and no length suffices if the bracket is not positive. With unlimited baselines
and $\sigma_o=0$, $T_{\mathrm{req}}=\zeta^{\star2}\phi S/(A\Gamma^2)$.

\subsection{Moments, signs and intervals}

The denominator $b=\hat\Gamma$ is a difference of two estimates. Under a normal
approximation it has positive density at zero, so $\mathbb{E}|\hat h|=\infty$ at every
$\nu$ (Hinkley, 1969; Marsaglia, 1965); with count data, $b=0$ has positive probability and
$\hat h$ is undefined there. No value of $\nu$ makes the moments exist.
What $\nu$ controls is the probability that $\hat\Gamma$ has the wrong sign,
$\Phi(-\nu)$: 0.023 at $\nu=2$ and 0.0013 at $\nu=3$.

The Fieller set for $h=a/b$ at level $1-\alpha$ is
\[
\{h:(a-hb)^2\le z_{\alpha/2}^2(v_a+h^2v_b)\},
\]
with $v_a,v_b$ the estimated variances of $a$ and $b$. It is a bounded interval
exactly when $b^2>z_{\alpha/2}^2v_b$, that is when $\hat\nu>1.96$ at the 95\% level.
Otherwise it is unbounded, which correctly reports that the dyad's data do not
bound $h$. In simulation, 95\% Fieller sets covered at 0.94 to 0.97 in every cell
without occasion variance.

\subsection{Why truncation and filtering bias toward coupling}

Under $\kappa=0$, $\hat h$ is centred near 1 and, at low $\zeta$, spreads widely to
both sides. Retaining only $0<\hat h<1$ discards the upper half of the null
distribution and every negative value. Every retained estimate maps to
$\hat\theta>0$, so the retained median is positive however large the sample. The
remedy is to invert over the whole real line, using $h$ extended to $\theta<0$ and
$\hat\theta=+\infty$ for $\hat h\le0$, and to summarise with order statistics or
Fieller sets.

Retaining only dyads with $\hat\nu\ge\nu_0$ conditions on $|\hat\Gamma|$ being large.
For a dyad of fixed true contrast near the threshold, this selects replicates in
which $|\hat\Gamma|>|\Gamma|$. The retained $\hat h=a/b$ then has an inflated
denominator and is biased toward 0, which reads as coupling: regression to the
mean (Barnett et al., 2005). The bias disappears if selection
uses sessions independent of those entering $\hat\Gamma$.

% =====================================================================
\section{Sampling variability of the gap}\label{sec:gap}

\paragraph{Dispersion.} For a counting process with dispersion $\phi$,
$\mathrm{Var}(k(T))\approx\phi\,\mathbb{E}k(T)$. For a stationary renewal process,
$\phi$ tends to the squared coefficient of variation of the inter-event intervals
(Cox, 1962); $\phi<1$ for processes more regular than Poisson, such as bites
separated by chewing. Independent counts for the two diners give
$\mathrm{Var}(D)=\phi(N_i+N_j)$, which at the FIC mean of 63.4 bites per meal is
$11.3^2$ for $\phi=1$ and $9.4^2$ for $\phi=0.7$. These are benchmarks, not floors.

\paragraph{Occasion variance.} If a person's rate on a given occasion is
$\beta+\varepsilon$ with $\varepsilon\sim N(0,\sigma_o^2)$ independent across
occasions, the slow-drive gap is $D=AhT(\Gamma+\varepsilon_j-\varepsilon_i)$ plus
counting noise. Then
\[
\sigma_D^2=\phi AST+2A^2h^2\sigma_o^2T^2 ,
\]
and the gap signal-to-noise ratio at $h=1$ is bounded by $\Gamma/(\sqrt2\,\sigma_o)$
for every $T$. The baseline contrast from $m$ sessions of length $T_s$ has
variance $(\phi S/T_s+2\sigma_o^2)/m$.

\paragraph{Coupling on realised events.} The simulations draw each diner's counts
independently given the deterministic means of main-text Eq.~(9). Suppose instead
that the coupling acts on realised counts, so that in a stochastic version of
Eq.~(1) the gap obeys, approximately,
\[
dD=(\gt-2\kappa D)\,dt+dM,
\]
where $M$ is a martingale with quadratic variation $\phi\lambda t$ and $\lambda$ is
the combined bite rate. In the slow-drive limit this is an Ornstein--Uhlenbeck
process, and
\[
\mathrm{Var}\,D(T)=\phi\lambda\int_0^Te^{-4\kappa(T-s)}\,ds=\phi\lambda T\,h(4\kappa T)=\phi\lambda T\,h(2\theta),
\]
which is smaller than $\phi\lambda T$ whenever $\kappa>0$. This is a diffusion
approximation and is not used in any result. It implies that the per-dyad
requirements are conservative for coupled dyads and exact under the null. It also
implies that the variance of the gap carries information about $\kappa$ even when
$\Gamma=0$, and that the two couplings may differ in this second moment.

% =====================================================================
\section{The pooled estimator}\label{sec:pooled}

\subsection{Model}

In the slow-drive limit of main-text Eq.~(9), each diner's intake is an
actor--partner combination of the two baselines, with actor effect $(1+h)/2$ and
partner effect $(1-h)/2$. For dyad $d$, with arbitrary labelling of $i$ and $j$,
define
\[
y_d=\hat S_d\,\frac{\kj-\ki}{\kj+\ki},\qquad
x_d=\hat\Gamma_d=\hat\beta_j-\hat\beta_i,\qquad
v_d=\frac{s_i^2+s_j^2}{m},
\]
where $\hat\beta$ is the mean of a person's $m$ per-session rates and $s^2$ their
sample variance. Then, to first order,
\[
y_d=h\Gamma_d+e_d,\qquad x_d=\Gamma_d+u_d,\qquad \mathbb{E}\,v_d=\mathrm{Var}(u_d),
\]
with $e_d$, $u_d$ and $\Gamma_d$ mutually independent and
\[
\sigma_e^2=\frac{\phi S}{AT}+2h^2\sigma_o^2,\qquad
\sigma_u^2=\frac{1}{m}\Bigl(\frac{\phi S}{T_s}+2\sigma_o^2\Bigr).
\]
$A$ and $T$ cancel from $y_d$, so neither needs to be estimated. Unbiasedness of
$v_d$ follows from $\mathbb{E}s^2=\sigma_o^2+\phi\beta/T_s$ for each person.

\subsection{Estimator and its variance}

The method-of-moments estimator solves $\sum_d\psi_d(h)=0$ with
$\psi_d(h)=y_dx_d-h(x_d^2-v_d)$:
\[
\hat h_{\mathrm{pool}}=\frac{\sum_dy_dx_d}{\sum_d(x_d^2-v_d)} .
\]
$\mathbb{E}\psi_d(h)=h\,\mathbb{E}\Gamma^2-h(\mathbb{E}\Gamma^2+\sigma_u^2-\sigma_u^2)=0$,
so the estimator is consistent. Its sandwich variance estimate is
$\sum_d\hat\psi_d^2/\bigl(\sum_d(x_d^2-v_d)\bigr)^2$ (Stefanski \& Boos, 2002). In small
samples, replace $\hat\psi_d$ by $\hat\psi_d/(1-\ell_d)$ with leverage
$\ell_d=(x_d^2-v_d)/\sum_d(x_d^2-v_d)$, the analogue of HC3 for least squares
(MacKinnon \& White, 1985); Section~\ref{sec:robust} shows that this restores
nominal size down to 20 dyads per group.

For the large-sample variance, write $\overline{\Gamma^2}=\mathbb{E}\Gamma^2$ and take
$\Gamma$ symmetric (random labels), $u$ normal, and $v$ independent of $u$, which
holds for normal per-session rates. Expanding,
\begin{align*}
\mathrm{Var}(y x)&=h^2\mathrm{Var}(\Gamma^2)+h^2\overline{\Gamma^2}\sigma_u^2+\sigma_e^2(\overline{\Gamma^2}+\sigma_u^2),\\
\mathrm{Var}(x^2-v)&=\mathrm{Var}(\Gamma^2)+4\overline{\Gamma^2}\sigma_u^2+2\sigma_u^4+\mathrm{Var}(v),\\
\mathrm{Cov}(yx,\,x^2-v)&=h\,\mathrm{Var}(\Gamma^2)+2h\overline{\Gamma^2}\sigma_u^2,
\end{align*}
with $\mathrm{Var}(v)=\sigma_u^4/(m-1)$ for equal per-person variances. Then
$\mathrm{Var}\,\psi=\mathrm{Var}(yx)+h^2\mathrm{Var}(x^2-v)-2h\,\mathrm{Cov}(yx,x^2-v)$, in
which $\mathrm{Var}(\Gamma^2)$ cancels:
\[
\mathrm{Var}\,\psi=\sigma_e^2(\overline{\Gamma^2}+\sigma_u^2)+h^2\overline{\Gamma^2}\sigma_u^2
+h^2\sigma_u^4\Bigl(2+\frac{1}{m-1}\Bigr).
\]
With $\mathbb{E}\,\partial\psi/\partial h=-\overline{\Gamma^2}$, the large-sample
variance over $n$ dyads is
\begin{equation}\label{eq:pooledvar}
\mathrm{Var}(\hat h_{\mathrm{pool}})\approx\frac{\mathrm{Var}\,\psi}{n\,(\overline{\Gamma^2})^2},
\end{equation}
main-text Eq.~(11). Heterogeneity in contrasts across dyads does not inflate it;
only their mean square enters. For a two-group comparison the test statistic is
$(\hat h_1-\hat h_2)/\sqrt{\widehat{\mathrm{se}}_1^2+\widehat{\mathrm{se}}_2^2}$, and
power at difference $\Delta h$ is
$\Phi(|\Delta h|/\mathrm{SE}-z_{\alpha/2})+\Phi(-|\Delta h|/\mathrm{SE}-z_{\alpha/2})$.
In simulation \eqref{eq:pooledvar} predicted power to within 0.032 for $m\ge4$ and
$n\ge50$ per group.

If coupling varies across dyads, $\hat h_{\mathrm{pool}}$ converges to
$\mathbb{E}[h_d\Gamma_d^2]/\mathbb{E}[\Gamma_d^2]$, a contrast-weighted mean.

\begin{remark}
The derivation assumes the slow-drive model, normal per-session rates and
independent counts, and the simulations reported in the main text use the same
assumptions. Their agreement with \eqref{eq:pooledvar} checks the large-sample
approximation under the model. Section~\ref{sec:robust} examines small samples,
skewed rates, a decaying drive, heterogeneous coupling and differential facilitation,
and Section~\ref{sec:facil} facilitation that differs between groups and pre-existing
pairs; coupling on realised events is not examined.
\end{remark}

\subsection{Pairing schemes}

For random pairing from a population with between-person standard deviation
$\sigma_p$, $\Gamma\sim N(0,2\sigma_p^2)$, so $\overline{\Gamma^2}=2\sigma_p^2$ and the
median $|\Gamma|$ is $0.954\,\sigma_p$.

For pairing opposite terciles of a screening session with reliability
$R=\sigma_p^2/(\sigma_p^2+\sigma_o^2+\phi\mu/T_{\mathrm{scr}})$, where $\mu$ is the
mean rate, write $\beta-\mu=R(X-\mu)+w$ with $X$ the screening rate and
$w\sim N(0,\sigma_p^2(1-R))$ independent of $X$. For a standard normal $Z$ above
its upper tercile point, $\mathbb{E}Z=k=1.0909$ and $\mathrm{Var}Z=v_t=0.2800$. Hence
$\mathbb{E}[\beta-\mu\mid\text{top}]=k\sqrt R\,\sigma_p$ and
$\mathrm{Var}[\beta\mid\text{top}]=\sigma_p^2(1-(1-v_t)R)$, and symmetrically for the
bottom tercile. So the mean contrast is $2k\sqrt R\,\sigma_p=2.18\sqrt R\,\sigma_p$
and
\[
\overline{\Gamma^2}=2\sigma_p^2\bigl(1-(1-v_t)R\bigr)+4k^2R\,\sigma_p^2
=\sigma_p^2(2+3.32R).
\]
The screening session must be independent of the baseline sessions: selecting on
the baselines themselves makes $u_d$ correlated with selection and biases the
estimator, as in Section~\ref{sec:perdyad}.

\subsection{Alternatives}

\paragraph{Instrumental variables.} Split each person's solo sessions into halves
with contrasts $x^{(1)}_d,x^{(2)}_d$, each with error variance $2\sigma_u^2$, and use
$\hat h_{\mathrm{IV}}=\sum y_dx^{(1)}_d/\sum x^{(2)}_dx^{(1)}_d$. With
$\psi=(y-hx^{(2)})x^{(1)}=(e-hu^{(2)})(\Gamma+u^{(1)})$,
\[
\mathrm{Var}(\hat h_{\mathrm{IV}})\approx\frac{(\sigma_e^2+2h^2\sigma_u^2)(\overline{\Gamma^2}+2\sigma_u^2)}{n(\overline{\Gamma^2})^2},
\]
which always exceeds \eqref{eq:pooledvar}. It does not require normality of the
per-session rates.

\paragraph{Naive least squares.}
$\hat h_{\mathrm{OLS}}=\sum y_dx_d/\sum x_d^2\to h\,\overline{\Gamma^2}/(\overline{\Gamma^2}+\sigma_u^2)$,
attenuated by the reliability of $\hat\Gamma$ (Fuller, 1987). Under $h=1$ this is
below 1, which reads as coupling. In simulation its null medians ranged from 0.50
to 0.97.

% =====================================================================
\section{Recommended protocol}

\begin{enumerate}
\item \emph{Size the study first.} Using pilot or published estimates of the mean
rate, the stable and occasion standard deviations of individual rates, and the
dispersion of the event process, compute the pooled sample size
(\texttt{pooled\_n}) and, if per-dyad estimates are of interest, the per-dyad
signal-to-noise ratio (\texttt{per\_dyad\_check}). Because these inputs are
uncertain, evaluate the sample size across a plausible range of them (main-text
Table~3C) and at the contrast in $h$ expected between groups, not only at
$\theta=0$ versus $\theta=1$. If the stable and occasion components are not known,
measure them first: with half the variation occasional, about 50 people with four
solo meals each estimate the occasion share to within about 0.2
(Section~\ref{sec:pilot}).
\item \emph{Solo sessions.} Record at least four solo sessions per participant
under conditions matching the joint session, and of the same length as the joint
session (Section~\ref{sec:robust}). Use the per-session rates for the
baseline (their mean) and for its error variance (their sample variance divided
by $m$). Report the stable and occasion variance components (a one-way
random-effects decomposition of per-session rates, net of counting variance).
\item \emph{Pairing.} Pair at random, or screen for contrasting rates using a
session that does not enter the baselines. If closeness groups consist of existing
pairs, which can be neither, measure the correlation of partners' solo rates and
use it in sizing (\texttt{pairing="existing"}).
\item \emph{Joint session.} Record event counts for both partners. $A$ and $T$ need
not be estimated, since both cancel.
\item \emph{Primary analysis.} Estimate $h$ per closeness group with
$\hat h_{\mathrm{pool}}$, with leverage-corrected sandwich or dyad-bootstrap
standard errors, and compare groups. Interpret the absolute level of
$\hat h_{\mathrm{pool}}$ only relative to the comparison group, since differential
facilitation can move it in either direction. Report the partner-effect test
(Section~\ref{sec:facil}) alongside: it does not need facilitation to be equal
across groups. A regression of $y_d$ on $x_d$ and $x_d\times$closeness, with the
same errors-in-variables correction, generalises this to a continuous moderator.
\item \emph{Per-dyad estimates, if reported.} Report them descriptively with
Fieller intervals and with $\zeta$. Do not truncate at $\hat h=1$, and do not
select dyads on the baselines that enter the estimate.
\end{enumerate}

% =====================================================================
\section{Worked example with \texttt{couplingdx}}

The package has no dependencies. It installs from the repository with
\begin{verbatim}
pip install "git+https://github.com/markizgarshev90-ops/couplingdx#subdirectory=couplingdx"
\end{verbatim}
or with \texttt{pip install .} from its folder. The values below are those of the
main text: mean bite rate $0.0867\,\mathrm{s}^{-1}$ per person, between-meal
standard deviation $0.0130\,\mathrm{s}^{-1}$, dispersion 0.7, 12-minute (740\,s)
meals.

\paragraph{Can a single dyad be estimated?}
\begin{verbatim}
>>> import couplingdx as c
>>> G = c.expected_contrast(0.0130)          # median contrast, random pairs
>>> round(G, 4)
0.0124
>>> print(c.per_dyad_check(contrast=G, rate_total=2*0.0867, duration=740,
...                        n_baseline=8, baseline_session=740, dispersion=0.7))
Per-dyad coupling estimate
========================================
  gap signal-to-noise      z0   =    0.97
  baseline contrast ratio  nu   =    2.74   (P sign error = 0.003)
  combined                 zeta =    0.94   (need >= 5.2)

  NOT FEASIBLE per dyad: theta_hat would be less precise than theta itself.
  required joint-session length: none suffices

Notes:
  - baseline precision alone rules out the target: more or longer
    baseline sessions are needed before joint-session length matters.
  - never discard estimates with h_hat >= 1, and never select dyads on
    the same baselines used in the estimate: both bias toward coupling.
\end{verbatim}

\paragraph{How long would the meal have to be?}
\begin{verbatim}
>>> c.required_duration(G, 2*0.0867, 0.7) / 3600          # unlimited baselines
5.93
>>> Gs = c.expected_contrast(0.0130, "screened", mean_rate=0.0867,
...                          dispersion=0.7, screen_duration=740)
>>> c.required_duration(Gs, 2*0.0867, 0.7, n_baseline=8,
...                     baseline_session=740) / 60             # minutes
171
\end{verbatim}

\paragraph{How many dyads does the pooled estimator need?}
\begin{verbatim}
>>> d = dict(duration=740, n_baseline=4, baseline_session=740,
...          mean_rate=0.0867, between_sd=0.0130, dispersion=0.7)
>>> c.pooled_n(0.8, **d)                                   # theta 0 vs 1
76
>>> c.pooled_n(0.8, **dict(d, between_sd=0.0130/2**0.5, occasion_sd=0.0130/2**0.5))
376
>>> c.pooled_n(0.8, **dict(d, pairing="screened"))
34
>>> c.pooled_n(0.8, **d,              # close group made of existing pairs
...            high_design=dict(pairing="existing", partner_correlation=0.3))
94
\end{verbatim}

\paragraph{From the command line.}
\begin{verbatim}
$ couplingdx pooled --mean-rate 0.0867 --between-sd 0.013 --duration 740 \
      --n-baseline 4 --baseline-session 740 --dispersion 0.7
Pooled estimator, theta 0 vs 1, two-sided 5% test
  dyads per group for power 0.80: 76
  power at   50 per group: 0.624
  power at  100 per group: 0.896
  power at  200 per group: 0.995
  power at  400 per group: 1.000
\end{verbatim}

\paragraph{Fitting the estimator.} Each dyad is given as its two members'
per-session solo rates and their joint-session counts, and each closeness group is
fitted separately:
\begin{verbatim}
>>> dyad = ([0.081, 0.090, 0.077, 0.085],   # solo rates, person i (bites/s)
...         [0.102, 0.095, 0.110, 0.099],   # solo rates, person j
...         66, 73)                          # joint-session counts, i and j
>>> fit_close = c.fit_pooled(dyads_close)        # list of such tuples
>>> fit_strangers = c.fit_pooled(dyads_strangers)
>>> c.compare_groups(fit_close, fit_strangers)   # diff, se, z, p
\end{verbatim}
The script \texttt{analysis/pooled\_fit.py} in the repository simulates 300 dyads
at $\theta=1$ and prints
\begin{verbatim}
h_pool = 0.612 (SE 0.047; 95% CI 0.519 to 0.704), theta = 1.08, n = 300 dyads
\end{verbatim}
\texttt{fit\_pooled} uses the leverage-corrected standard error by default and
refuses to return an estimate if the corrected denominator is not clearly positive.
The partner-effect test of Section~\ref{sec:facil} takes the same dyads and the
joint-session length:
\begin{verbatim}
>>> print(c.fit_partner_effect(dyads_close, duration=740))
\end{verbatim}

Below about 50 dyads per group the large-sample formula overstates power by up to
0.07 (Section~\ref{sec:robust}); such designs, and designs with departures not
covered there, should be checked by simulation, for which
\texttt{robustness\_sim.py} can be adapted.

% =====================================================================
\section{Robustness of the pooled estimator}\label{sec:robust}

The simulations of main-text Figure~4 generate data under the assumptions of
Eq.~\eqref{eq:pooledvar}. Table~\ref{tab:robust} reports a second set that departs
from them one at a time. Unless stated, each cell has two groups of 100 dyads
($\theta=0$ and $\theta=1$), 12-minute joint and solo sessions ($T=T_s=740$\,s), $m=4$
solo sessions, random pairing, $\phi=0.7$, mean rate $0.0867\,\mathrm{s}^{-1}$ and
$\sigma_b=0.0130\,\mathrm{s}^{-1}$, all stable; 4{,}000 replicates per cell; and the
leverage-corrected (HC3) standard error of main-text Method. The code is
\texttt{robustness\_sim.py} in the repository.

\begin{description}\itemsep2pt
\item[R1, small samples.] 20 or 30 dyads per group, random or screened pairing. With
the plain sandwich (HC0) the null rejection rates were 0.087 at 20 and 0.073 at 30
dyads per group (6{,}000 replicates; \texttt{hc\_comparison.py}); HC3 divides each score by $1-\ell_d$, $\ell_d=(x_d^2-v_d)/\sum_d(x_d^2-v_d)$.
\item[R2, skewed rates.] Stable rates lognormal with the same mean and standard
deviation; occasion effects a centred gamma variable with shape 2 (skewness 1.4).
\item[R3, decaying drive.] All sessions generated from the exact solution of main-text
Eq.~(1), Eq.~\eqref{eq:solution}, with $l_p\sim U(0,l_{\max})$ independently per
person; solo rates computed as count divided by session length, as an analyst would.
\item[R4, heterogeneous coupling.] $\theta_d\sim\mathrm{Gamma}(\text{shape }2,
\text{scale }0.5)$ in the coupled group. The target is the contrast-weighted mean
$\mathbb{E}[h_d\Gamma_d^2]/\mathbb{E}[\Gamma_d^2]$ (Section~\ref{sec:pooled}).
\item[R5, differential facilitation.] $A_p=1.1\,(1+\delta z_p)$, with $z_p$ the
person's standardised solo rate, applied in both groups.
\end{description}

\begin{table}[H]
\centering\footnotesize
\caption{Robustness of the pooled estimator. Medians of $\hat h_{\mathrm{pool}}$ in
the uncoupled ($\theta=0$) and coupled ($\theta=1$) groups, with the target in
parentheses; coverage of $h=1$ by 95\% intervals in the uncoupled group; null
rejection rate of the two-group test (size); power; and the power predicted by
main-text Eq.~(11). Monte Carlo standard errors are about 0.003 for
size and at most 0.008 for power.}\label{tab:robust}
\setlength{\tabcolsep}{3.5pt}
\begin{tabular}{@{}lrrrrrrr@{}}
\toprule
& & \multicolumn{2}{c}{Uncoupled, $\theta=0$} & Coupled, $\theta=1$ & \multicolumn{3}{c}{Two-group test}\\
\cmidrule(lr){3-4}\cmidrule(lr){5-5}\cmidrule(l){6-8}
Scenario & $n$ & Med.\ $\hat h$ & Cover & Med.\ $\hat h$ (target) & Size & Power & Formula\\
\midrule
Model as assumed & 100 & 1.000 & 0.944 & 0.633 (0.632) & 0.051 & 0.886 & 0.896\\
R1: 20 dyads per group & 20 & 1.004 & 0.940 & 0.636 (0.632) & 0.050 & 0.255 & 0.302\\
R1: 30 dyads per group & 30 & 1.006 & 0.936 & 0.636 (0.632) & 0.051 & 0.373 & 0.422\\
R1: 20 dyads, screened & 20 & 1.001 & 0.932 & 0.639 (0.632) & 0.050 & 0.513 & 0.583\\
R2: normal, half occasional & 300 & 1.012 & 0.953 & 0.637 (0.632) & --- & 0.696 & 0.707\\
R2: skewed, half occasional & 300 & 1.009 & 0.955 & 0.636 (0.632) & 0.046 & 0.718 & 0.707\\
R2: skewed, all stable & 100 & 1.001 & 0.951 & 0.633 (0.632) & --- & 0.868 & 0.896\\
R3: $l_{\max}=2.4\times10^{-4}$ & 100 & 1.000 & 0.943 & 0.632 (0.632) & 0.048 & 0.866 & 0.896\\
R3: $l_{\max}=5\times10^{-4}$, $T=T_s=740$ & 100 & 1.002 & 0.951 & 0.638 (0.632) & --- & 0.883 & 0.896\\
R3: $l_{\max}=5\times10^{-4}$, $T=T_s=3000$ & 100 & 1.001 & 0.941 & 0.668 (0.632) & --- & 1.000 & 1.000\\
R3: $l_{\max}=5\times10^{-4}$, $T=3000$, $T_s=740$ & 100 & 1.269 & 0.266 & 0.803 (0.632) & 0.045 & 0.933 & 1.000\\
R3: no decay, $T=3000$, $T_s=740$ & 100 & 1.003 & 0.945 & 0.633 (0.632) & --- & 0.999 & 1.000\\
R4: $\theta_d\sim\mathrm{Gamma}(2,0.5)$ & 100 & 1.001 & 0.943 & 0.669 (0.667) & --- & 0.792 & 0.896\\
R5: $\delta=-0.05$ & 100 & 0.674 & 0.024 & 0.427 (0.632) & 0.045 & 0.621 & 0.896\\
R5: $\delta=+0.05$ & 100 & 1.322 & 0.032 & 0.834 (0.632) & 0.054 & 0.980 & 0.896\\
\bottomrule
\end{tabular}
\end{table}

Three results matter for design. First, HC3 is needed below about 50 dyads per
group, and the formula then overstates power by up to 0.07. Second, decay is
harmless at the rates FIC allows and remains harmless under the null at twice
those rates, provided solo and joint sessions have equal length. With unequal
lengths it biases $\hat h_{\mathrm{pool}}$ strongly, although the between-group test
keeps its size. Third, differential facilitation and heterogeneous coupling bias
$\hat h_{\mathrm{pool}}$, or reduce its power, without affecting the size of a
comparison between groups that share them. The absolute level of
$\hat h_{\mathrm{pool}}$ should therefore be interpreted only relative to a
comparison condition. Section~\ref{sec:facil} examines facilitation that differs
between groups, and groups made of pre-existing pairs.


% =====================================================================
\section{Facilitation that differs between groups, the partner-effect test, and pre-existing pairs}\label{sec:facil}

\paragraph{Design.} Each cell has 100 dyads per group and 4{,}000 replicates, with the
settings of Section~\ref{sec:robust} and $A_p=1.1\,(1+\delta z_p)$. In R6 the
strangers have $\delta=0$ and the close group $\delta=\pm0.05$. In R7 the close group
consists of pre-existing pairs whose stable rates are bivariate normal with
correlation $\rho$, and the strangers are paired at random. In the
contrast-dependent case both groups follow the same rule,
$\theta_d=2\exp(-|\Gamma_d|/\mathbb{E}|\Gamma|)$ with $\mathbb{E}|\Gamma|$ taken over
random pairs, so that more similar pairs couple more. The code is
\texttt{revision\_sims.py} in the repository.

\paragraph{The partner-effect test.} Without the constraint that the actor and
partner effects sum to one, the slow-drive intakes are
\begin{equation}\label{eq:apimfree}
\frac{\ki(T)}{T}=\frac{1+h}{2}A_i\beta_i+\frac{1-h}{2}A_j\beta_j ,
\end{equation}
with $A_p$ allowed to depend on the person. Each diner's joint-session rate is
regressed on their own and their partner's baseline, with an intercept, by the
method of moments: the summed replicate variances $\sum v$ are subtracted from the
diagonal of the moment matrix of the centred baselines, as in
Section~\ref{sec:pooled}. Both members of a dyad enter as rows, and the standard
error of the partner coefficient comes from deleting one dyad at a time.

Under $\kappa=0$ a diner's joint-session rate depends on their own rate only,
$A_i\beta_i=f(\beta_i)$ for some function $f$. Its coefficient on the partner's rate,
in a linear regression on both rates, is zero whenever $\mathbb{E}[\beta_j\mid\beta_i]$
is linear in $\beta_i$, because the part of $f(\beta_i)$ that is not linear in
$\beta_i$ is then uncorrelated with $\beta_j$. This holds for random pairing, where
$\beta_j$ is independent of $\beta_i$, and for pre-existing pairs with jointly
normal rates at any correlation. Under coupling the partner coefficient is
$(1-h)/2$ times the slope of $A_j\beta_j$ in $\beta_j$, which is $A(1-h)/2$ when $A$
is constant. A gradient $\delta$ multiplies that slope by about
$1+\delta\mu/\sigma_p$ near the mean rate, 1.33 for $\delta=0.05$ at the FIC
values, which is why the power of the test depends on $\delta$
(Table~\ref{tab:facil}).

\begin{table}[H]
\centering\footnotesize
\caption{Facilitation that differs between groups (R6), 100 dyads per group.
A: rejection rates of the between-group test of $h$ (pooled estimator), of the
between-group difference in partner effects, and of the partner-effect test in
the close group alone, with no coupling in either group ($\theta=0$) and with
$\theta=1$ in the close group. B: the partner-effect test in one group, its median
estimate, its target $A(1-h)/2$ with $A=1.1$, and its rejection rate.}\label{tab:facil}
\setlength{\tabcolsep}{4pt}
\begin{tabular}{@{}lrrrrrr@{}}
\toprule
& \multicolumn{2}{c}{Pooled $h$, groups} & \multicolumn{2}{c}{Partner, groups} & \multicolumn{2}{c}{Partner, close}\\
\cmidrule(lr){2-3}\cmidrule(lr){4-5}\cmidrule(l){6-7}
A. Scenario & $\theta=0$ & $\theta=1$ & $\theta=0$ & $\theta=1$ & $\theta=0$ & $\theta=1$\\
\midrule
Close group only, $\delta=-0.05$ & 0.826 & 0.999 & 0.045 & 0.349 & 0.050 & 0.638\\
Close group only, $\delta=+0.05$ & 0.749 & 0.324 & 0.052 & 0.827 & 0.055 & 0.984\\
\bottomrule
\end{tabular}

\vspace{6pt}
\begin{tabular}{@{}lrrr@{}}
\toprule
B. Scenario & Median & Target & Rejection\\
\midrule
Random pairing, $\delta=-0.05$, $\theta=0$ & 0.002 & --- & 0.055\\
Random pairing, $\delta=+0.00$, $\theta=0$ & 0.000 & --- & 0.052\\
Random pairing, $\delta=+0.05$, $\theta=0$ & 0.002 & --- & 0.052\\
Screened pairing, $\delta=-0.05$, $\theta=0$ & 0.002 & --- & 0.050\\
Screened pairing, $\delta=+0.00$, $\theta=0$ & 0.002 & --- & 0.045\\
Screened pairing, $\delta=+0.05$, $\theta=0$ & 0.002 & --- & 0.047\\
Random pairing, $\delta=0$, $\theta=0.5$ & 0.118 & 0.117 & 0.474\\
Random pairing, $\delta=0$, $\theta=1$ & 0.204 & 0.202 & 0.907\\
\bottomrule
\end{tabular}
\end{table}

\begin{table}[H]
\centering\footnotesize
\caption{Closeness groups made of pre-existing pairs (R7), 100 dyads per group;
strangers are paired at random. Rejection rate of the between-group test with no
coupling in either group (or, for coupling that falls with contrast, with the same
rule in both groups), its power at $\theta=1$ in the close group, the power from
main-text Eq.~(11) with mean squared contrast $2\sigma_p^2(1-\rho)$
in the close group, and the median estimate in the close group: $\hat h_{\mathrm{pool}}$
(strangers in parentheses), or the partner effect for the partner-effect test,
which has $\theta=0$.}\label{tab:existing}
\setlength{\tabcolsep}{4pt}
\begin{tabular}{@{}lrrrr@{}}
\toprule
Scenario & Rejection, null & Power & Formula & Median estimate\\
\midrule
Constant coupling, $\rho=0.3$ & 0.060 & 0.791 & 0.826 & 0.632\\
Constant coupling, $\rho=0.6$ & 0.052 & 0.601 & 0.645 & 0.635\\
Coupling falls with contrast, $\rho=0.3$ & 0.058 & --- & --- & 0.813 (0.848)\\
Coupling falls with contrast, $\rho=0.6$ & 0.088 & --- & --- & 0.748 (0.847)\\
Partner-effect test, $\rho=0.3$, $\delta=0.00$ & 0.049 & --- & --- & 0.002\\
Partner-effect test, $\rho=0.3$, $\delta=0.05$ & 0.055 & --- & --- & -0.001\\
Partner-effect test, $\rho=0.6$, $\delta=0.00$ & 0.052 & --- & --- & -0.003\\
Partner-effect test, $\rho=0.6$, $\delta=0.05$ & 0.053 & --- & --- & -0.001\\
\bottomrule
\end{tabular}
\end{table}

With constant coupling, correlated partners cost power but not size. When coupling
falls with the pair's contrast, the close group's more similar pairs couple more on
average ($\bar\theta=1.02$ at $\rho=0.3$ and $1.17$ at $\rho=0.6$, against $0.92$
among strangers), so a difference between the groups appears that closeness did not
cause.

% =====================================================================
\section{Group-level tests on per-dyad estimates}\label{sec:rank}

Two groups of randomly paired dyads were simulated with 12-minute joint and solo
sessions, $m=4$, $\phi=0.7$, and all between-meal variation stable
($\sigma_p=0.0130\,\mathrm{s}^{-1}$) or half occasional, 2{,}000 replicates per cell.
Per-dyad estimates from Eq.~(\ref{eq:estimator}) were not truncated; a dyad whose
estimated contrast was exactly zero, which is rare, was given $\hat h=0$. The groups
were compared with a two-sided Mann--Whitney test and, on the same data, with the
pooled estimator.

\begin{table}[H]
\centering\footnotesize
\caption{Rank test on untruncated per-dyad estimates against the pooled estimator:
rejection rates with $\theta=0$ in both groups (size) and $\theta=0$ against
$\theta=1$ (power), and the median per-dyad $\hat h$ in a group with no coupling.}\label{tab:rank}
\begin{tabular}{@{}lrrrrrr@{}}
\toprule
& & \multicolumn{2}{c}{Rank test} & \multicolumn{2}{c}{Pooled} & \\
\cmidrule(lr){3-4}\cmidrule(lr){5-6}
Variance & Dyads & Size & Power & Size & Power & Median $\hat h$, $\theta=0$\\
\midrule
All stable & 50 & 0.051 & 0.229 & 0.051 & 0.571 & 0.89\\
 & 100 & 0.057 & 0.422 & 0.048 & 0.869 & 0.89\\
 & 200 & 0.041 & 0.705 & 0.053 & 0.990 & 0.88\\
 & 400 & 0.041 & 0.944 & 0.042 & 1.000 & 0.89\\
Half occasional & 50 & 0.049 & 0.101 & 0.035 & 0.133 & 0.65\\
 & 100 & 0.061 & 0.148 & 0.041 & 0.292 & 0.66\\
 & 200 & 0.050 & 0.233 & 0.039 & 0.509 & 0.66\\
 & 400 & 0.049 & 0.398 & 0.043 & 0.806 & 0.66\\
\bottomrule
\end{tabular}
\end{table}

The rank test held its size. At the power levels that matter it needed about three
times as many dyads as the pooled estimator: the power it reached with 400 dyads per
group, the pooled estimator reached with about 130. Under no coupling the median
per-dyad estimate was well below 1. For a dyad whose true contrast is small,
$\hat h$ is close to a ratio of two independent noise terms, whose median is zero,
and such dyads pull the median of a randomly paired population below 1.

% =====================================================================
\section{What the time course adds}\label{sec:timecourse}

In the slow-drive limit, with $\Gamma_o$ the contrast on the joint occasion
(including occasion shifts and facilitation),
$D(T/2)=\Gamma_o(T/2)\,h(\theta/2)$ and
\begin{equation}\label{eq:halves}
D(T)-D(T/2)=e^{-\theta/2}D(T/2),\qquad \frac{D(T)}{D(T/2)}=1+e^{-\theta/2}.
\end{equation}
Let $D_1$ be the gap at half-time, $D_2$ the change in the gap over the second
half, and $N_1$ the two diners' bites in the first half. If the counts in the two
halves are independent given the drive, with variance $\phi$ times their mean, then
$\mathbb{E}[D_1D_2]=e^{-\theta/2}\,\mathbb{E}[D_1^2-\phi N_1]$, and
\[
\hat\varrho=\frac{\sum_dD_{1d}D_{2d}}{\sum_d\bigl(D_{1d}^2-\phi N_{1d}\bigr)},\qquad
\hat\theta=-2\log\hat\varrho ,
\]
with a leverage-corrected sandwich standard error as in Section~\ref{sec:pooled}.
It needs no solo sessions, and an occasion shift in the joint session is part of
$\Gamma_o$ rather than noise. It does need $\phi$ for counts over half a meal.

The simulations follow Section~\ref{sec:robust}, with the counts of each half drawn
separately. The departures were: half-meal counts with dispersion 0.9 analysed with
0.7; a change in rate between the halves shared by both diners and varying across
dyads (mean $+4\%$, standard deviation 10\% on the log scale), in the joint session
only; faster eaters slowing by 5\% more per standard deviation of their rate in the
second half, in solo and joint sessions alike; and a decaying drive.

\begin{table}[H]
\centering\footnotesize
\caption{Pooled half-meal estimator against the end-of-meal estimator, 12-minute
meals, $m=4$ solo meals for the end-of-meal estimator. A: median $\hat\varrho$ under
no coupling (target 1) and at $\theta=1$ (target 0.607), and rejection rates of
both two-group tests. B: departures, 100 dyads per group; median estimate under no
coupling and rejection rates.}\label{tab:timecourse}
\setlength{\tabcolsep}{4pt}
\begin{tabular}{@{}lrrrrrrr@{}}
\toprule
& & \multicolumn{4}{c}{Half-meal estimator} & \multicolumn{2}{c}{End-of-meal}\\
\cmidrule(lr){3-6}\cmidrule(l){7-8}
A. Variance & Dyads & $\hat\varrho$, $\theta=0$ & $\hat\varrho$, $\theta=1$ & Size & Power & Size & Power\\
\midrule
All stable & 50 & 1.012 & 0.617 & 0.007 & 0.065 & 0.048 & 0.593\\
 & 100 & 1.005 & 0.619 & 0.019 & 0.167 & 0.051 & 0.885\\
 & 200 & 1.003 & 0.605 & 0.033 & 0.359 & 0.057 & 0.991\\
 & 400 & 0.996 & 0.607 & 0.041 & 0.615 & 0.046 & 1.000\\
Half occasional & 50 & 1.003 & 0.586 & 0.006 & 0.062 & 0.030 & 0.120\\
 & 100 & 1.004 & 0.622 & 0.014 & 0.169 & 0.054 & 0.261\\
 & 200 & 0.998 & 0.606 & 0.033 & 0.361 & 0.053 & 0.541\\
 & 400 & 0.999 & 0.605 & 0.037 & 0.633 & 0.046 & 0.824\\
\bottomrule
\end{tabular}

\vspace{6pt}
\begin{tabular}{@{}lrrrrrr@{}}
\toprule
& \multicolumn{3}{c}{Half-meal estimator} & \multicolumn{3}{c}{End-of-meal}\\
\cmidrule(lr){2-4}\cmidrule(l){5-7}
B. Departure & $\hat\varrho$, $\theta=0$ & Size & Power & $\hat h$, $\theta=0$ & Size & Power\\
\midrule
Half-meal dispersion 0.9, assumed 0.7 & 0.792 & 0.033 & 0.248 & 1.005 & 0.048 & 0.783\\
Shared second-half change, mean $+4\%$, SD 10\% & 1.046 & 0.017 & 0.178 & 1.002 & 0.050 & 0.875\\
Faster eaters slow 5\% more per SD & 0.666 & 0.025 & 0.242 & 1.003 & 0.057 & 0.798\\
Decaying drive, $l_{\max}=2.4\times10^{-4}$ & 0.993 & 0.018 & 0.170 & 1.006 & 0.054 & 0.875\\
Decaying drive, $l_{\max}=5\times10^{-4}$ & 0.975 & 0.017 & 0.185 & 1.001 & 0.057 & 0.888\\
\bottomrule
\end{tabular}
\end{table}

The half-meal test was conservative and much less powerful. The end-of-meal ratio is
unaffected by the shared change in rate, which cancels between the diners, and by
the person-specific slowing, which enters the solo and joint sessions alike; the
half-meal estimator is moved by both, and by an error in $\phi$. In FIC, the log
ratio of second-half to first-half bites had mean 0.037 and standard deviation
0.222 across meals, against 0.221 expected from counting noise at $\phi=0.7$.

% =====================================================================
\section{Pilot study size for the variance components}\label{sec:pilot}

$k$ people each ate $m$ solo 12-minute meals. Rates were normal with between-meal
standard deviation $0.0130\,\mathrm{s}^{-1}$, split as stated, counts had
$\phi=0.7$, and the components were estimated as in \texttt{variance\_components.py}:
the pooled within-person variance of the rates net of counting variance gives the
occasion part, and the between-meal variance net of counting variance gives the
total. Table~\ref{tab:pilot} gives the central 95\% range of the estimates over
2{,}000 replicates. With its actual design (five of 12 participants with repeat
meals), FIC gave an occasion share of 1, with a bootstrap interval from 0 to 1
(main text).

\begin{table}[H]
\centering\footnotesize
\caption{Central 95\% range of the estimated occasion share
$\sigma_o^2/(\sigma_o^2+\sigma_p^2)$ and stable standard deviation $\sigma_p$
($\mathrm{s}^{-1}$), true share in parentheses. FIC has 12 people and 21 meals.}\label{tab:pilot}
\begin{tabular}{@{}lrrrr@{}}
\toprule
Variance (true share) & People & Meals each & Occasion share & $\sigma_p$\\
\midrule
Half occasional (0.50) & 12 & 2 & 0.00--1.00 & 0.000--0.016\\
 & 12 & 4 & 0.12--0.99 & 0.001--0.014\\
 & 24 & 2 & 0.00--1.00 & 0.000--0.014\\
 & 24 & 4 & 0.23--0.83 & 0.005--0.013\\
 & 48 & 2 & 0.14--0.91 & 0.003--0.013\\
 & 48 & 4 & 0.30--0.72 & 0.006--0.012\\
 & 96 & 2 & 0.27--0.77 & 0.006--0.012\\
 & 96 & 4 & 0.36--0.67 & 0.007--0.011\\
All stable (0.00) & 12 & 2 & 0.00--0.62 & 0.005--0.019\\
 & 12 & 4 & 0.00--0.31 & 0.007--0.018\\
 & 24 & 2 & 0.00--0.35 & 0.008--0.017\\
 & 24 & 4 & 0.00--0.19 & 0.008--0.017\\
 & 48 & 2 & 0.00--0.23 & 0.009--0.016\\
 & 48 & 4 & 0.00--0.12 & 0.010--0.016\\
 & 96 & 2 & 0.00--0.15 & 0.010--0.015\\
 & 96 & 4 & 0.00--0.09 & 0.011--0.015\\
\bottomrule
\end{tabular}
\end{table}

\bigskip
\noindent\textbf{References}

\begingroup\small
\setlength{\parindent}{-0.3in}\setlength{\leftskip}{0.3in}\setlength{\parskip}{3pt}

Barnett, A. G., van der Pols, J. C., \& Dobson, A. J. (2005). Regression to the mean: What it is and how to deal with it. \emph{International Journal of Epidemiology, 34}(1), 215--220.

Cox, D. R. (1962). \emph{Renewal theory}. Methuen.

Fuller, W. A. (1987). \emph{Measurement error models}. Wiley.

Hinkley, D. V. (1969). On the ratio of two correlated normal random variables. \emph{Biometrika, 56}(3), 635--639.

Kieckens, S. N., \& Izgarshev, M. D. (2026, September 15--21). \emph{A dynamical model of social eating: Coupling, mimicry, and convergence in dyads} [Paper presentation]. International Conference of Numerical Analysis and Applied Mathematics (ICNAAM 2026), Heraklion, Greece.

MacKinnon, J. G., \& White, H. (1985). Some heteroskedasticity-consistent covariance matrix estimators with improved finite sample properties. \emph{Journal of Econometrics, 29}(3), 305--325.

Marsaglia, G. (1965). Ratios of normal variables and ratios of sums of uniform variables. \emph{Journal of the American Statistical Association, 60}(309), 193--204.

Stefanski, L. A., \& Boos, D. D. (2002). The calculus of M-estimation. \emph{The American Statistician, 56}(1), 29--38.

\endgroup

\end{document}
